\section{The normal form}\label{sec:normalform}

This section proves \cref{thm:nf-intro} and sets up the lattice on which all later
proofs live. The argument is short: relabelling acts on the sign configurations,
there are four orbits, and one operator identity on the invariant space collapses
each orbit to two triple products.

\subsection{Orbits of sign configurations}

Five tangent directions in general position, that is with no three of them coplanar,
define ten nonzero signs $\epsilon(IJK)$, $I<J<K\le5$; such a sign vector is a uniform
chirotope of rank three, so it satisfies the three-term Grassmann--Pl\"ucker relations.
This is the hypothesis of \cref{thm:nf-intro}; \cref{rem:nongeneric} says what
changes without it.

\begin{lemma}[Relabelling]\label{lem:relabel}
Let $\pi$ permute the edges. The unitary that permutes the tensor factors maps the
operator of $(\epsilon;j_1,\dots,j_5)$ to the operator of
$(\epsilon^\pi;j_{\pi(1)},\dots,j_{\pi(5)})$, where $\epsilon^\pi$ is the relabelled
chirotope. Moreover $A_{-\epsilon}=-A_\epsilon$.
\end{lemma}

\begin{proof}
$A_\epsilon=\frac16\sum\chi_\epsilon(I,J,K)\,\hat q_{IJK}$ over ordered triples of
distinct edges, where $\chi_\epsilon$ is the alternating extension of $\epsilon$; the
operator $\hat q_{IJK}=-4iJ_I\cdot(J_J\times J_K)$ is alternating because the three
vector operators commute. This expression is covariant under relabelling.
\end{proof}

So the set of spectra over all spin assignments is an invariant of the relabelling
orbit. The orbits are found by an exact finite computation (\cref{sec:computer}).

\begin{proposition}{Orbits}{orbits}
Exactly $384$ of the $2^{10}$ sign vectors satisfy the three-term
Grassmann--Pl\"ucker relations, and each is realised by integer tangent vectors with
entries in $[-4,4]$. Under relabelling they form four orbits, of sizes $24$, $120$,
$120$ and $120$, each closed under $\epsilon\mapsto-\epsilon$. Through
\eqref{eq:sigma} they produce $1$, $20$, $30$ and $120$ distinct reduced vectors
$\sigma$, $171$ in all ($85$ nonzero up to sign), and no $\sigma$ occurs in two orbits.
Every $\sigma$ of the last orbit contains $\pm4$, and no other $\sigma$ does.
\end{proposition}

These counts are Brunnemann and Rideout's: $384$ configurations and $171$ reduced
vectors at valence five~\cite[Table~1]{BR2008b}, and four classes under
permutation~\cite[Table~3]{BR2010}, whose representatives $(2,0,0,0)$, $(2,2,2,0)$ and
$(2,2,4,0)$ lie in the second, third and fourth orbit.

\subsection{Two triple products}

\begin{lemma}\label{lem:collapse}
On $\Inv(j_1\otimes\cdots\otimes j_5)$, $J_1\cdot\bigl((J_2+J_3)\times J_4\bigr)=
J_1\cdot(J_4\times J_5)$.
\end{lemma}

\begin{proof}
Write $B=J_2+J_3$. On invariant vectors $J_5=-(J_1+B+J_4)$, acting first. The terms
$J_1\cdot(J_4\times J_1)=-iJ_4\cdot J_1$ and $J_1\cdot(J_4\times J_4)=iJ_1\cdot J_4$
cancel, so $J_1\cdot(J_4\times J_5)=-J_1\cdot(J_4\times B)=J_1\cdot(B\times J_4)$.
\end{proof}

\begin{proof}[Proof of \cref{thm:nf-intro}]
For each orbit we exhibit a labelling with $\sigma=(\alpha,\beta,\beta,0)$: the orbits
give $(\alpha,\beta)=(2,0)$, $(2,2)$ and $(4,2)$ up to signs (the computation of
\cref{prop:orbits}). Then by \eqref{eq:sigma}
\[
  A=\alpha\hat q_{123}+\beta(\hat q_{124}+\hat q_{134})
   =-4i\bigl[\alpha J_1\cdot(J_2\times J_3)+\beta J_1\cdot\bigl((J_2+J_3)\times J_4\bigr)\bigr],
\]
which is \eqref{eq:nf} by \cref{lem:collapse}. The signs of $\alpha$ and $\beta$ do not
affect the spectrum: in the basis of the next subsection the first product preserves
$c=|J_4+J_5|$ and the second shifts it by $\pm1$, so the diagonal unitary $(-1)^c$
reverses the sign of $\beta$ alone, and $(-1)^b$ that of $\alpha$. Finally, $|\beta/\alpha|$
takes the value $\frac12$ only on the last orbit, which is the orbit of the $\sigma$
containing $\pm4$.
\end{proof}

For each $\epsilon$ we also solved the linear equation
$\sum\epsilon(IJK)\,J_I\cdot(J_J\times J_K)=\lambda J_a\cdot(J_b\times J_c)+\mu J_a\cdot
(J_d\times J_e)$ on the invariant space over all choices of the leg $a$ and of the
split of the other four. In the last two orbits exactly one leg $a$ admits a
solution with $\lambda\mu\ne0$; this is the \emph{distinguished edge}. In the third
orbit the split into pairs is not unique, because on the invariant space
\begin{equation}\label{eq:twosplits}
  J_1\cdot(J_2\times J_3+J_4\times J_5)=J_1\cdot(J_3\times J_4-J_2\times J_5),
\end{equation}
by the computation of \cref{lem:collapse} applied to
$J_1\cdot((J_2+J_4)\times(J_3+J_5))$. This is one face of a larger symmetry of the
$t=1$ operator, \cref{prop:S4} below.

\begin{remark}[Tangents not in general position]\label{rem:nongeneric}
If some three tangents are coplanar but not all five, some signs $\epsilon(IJK)$ vanish
and \cref{thm:nf-intro} does not apply. The same exact enumeration, run over
$\{0,\pm1\}^{10}$, finds $1\,680$ such sign vectors with pairwise non-parallel tangents
(each satisfies the chirotope axioms and is realised by integer tangents with entries in
$[-4,4]$), in $16$ relabelling orbits, and their reduced signs take every value from $-4$
to $4$. In $9$ orbits the operator still has the form \eqref{eq:nf}, now with
$|t|\in\{0,\frac13,\frac12,\frac23,1\}$ and $|\alpha|\in\{1,2,3\}$, so the results of
\cref{sec:closing,sec:tone} apply to it after rescaling $q$ by $|\alpha|/2$; in particular
\cref{thm:closing}(i) closes the gap on the four orbits with $|t|\in\{\frac13,\frac12,\frac23\}$.
In the other $7$ orbits ($780$ sign vectors) the operator is not of the form
\eqref{eq:nf}: in one of them it is a combination of two triple products that share two
edges, in the other six no two triple products suffice, and we do not treat them. Edges with parallel tangents are
the degenerate case of the remark at the end of \cref{sec:v6}.
\end{remark}

\subsection{The recoupling grid}

Fix the normal form \eqref{eq:nf} and write $s=j_1$ for the distinguished spin. The
basis adapted to it couples the two pairs,
\[
  |b,c\rangle,\qquad b=|J_2+J_3|,\quad c=|J_4+J_5|,\qquad |b-c|\le s\le b+c .
\]
For a pair $P=(j_a,j_b)$ put $n_P=j_a+j_b+1$ and $m_P=|j_a-j_b|$, so that the pair
spin ranges over $[m_P,n_P-1]$, and write $h_P(L)$ for the Stark links
\eqref{eq:stark} of $H(n_P,m_P)$. The sites $(b,c)$ fill a rectangle cut by the
strip $|d|\le s$ and the corner $u\ge s$, where
\[
  d=b-c,\qquad u=b+c,\qquad F(x)=(s+\tfrac12)^2-x^2,\qquad G(u)=(u+s+1)(u-s).
\]
The product $X=J_1\cdot(J_2\times J_3)$ commutes with $c$ and moves $b$ by one; on
fixed $c$ it is the four-valent operator of $\Inv(j_2,j_3,s,c)$. The product
$D=J_1\cdot(J_4\times J_5)$ commutes with $b$ and moves $c$.

\begin{proposition}{The grid}{grid}
A diagonal unitary gauge carries $iA_{\alpha,\beta}$, up to an overall sign, to the real
symmetric, entrywise nonnegative matrix on the sites with links
\begin{align*}
  (b-1,c)\text{--}(b,c):\quad &|\alpha|\,h_1(b)\sqrt{F(d-\tfrac12)\,G(u)},\\
  (b,c-1)\text{--}(b,c):\quad &|\beta|\,h_2(c)\sqrt{F(d+\tfrac12)\,G(u)},
\end{align*}
where $(b,c)$ is the upper end of the link and $d,u$ are taken there. In particular the
eigenvalue moduli of $A_{\alpha,\beta}$ are those of this matrix.
\end{proposition}

\begin{proof}
By \eqref{eq:ak} applied to $\Inv(j_2,j_3,s,c)$, the $X$-link has squared modulus
$\frac1{16}h_1(b)^2P(b;s,c)$ with $P(b;s,c)=((s+c+1)^2-b^2)(b^2-(s-c)^2)$, and the
polynomial identity
\begin{equation}\label{eq:heron}
  P(b;s,c)=F(d-\tfrac12)\,G(u),\qquad P(c;s,b)=F(d+\tfrac12)\,G(u)
\end{equation}
gives the stated form after multiplying by $4|\alpha|$; the $D$-links are the same with
the pairs exchanged. It remains to remove the signs. On invariant vectors
$X=-L\cdot(J_2\times J_3)$ with $L=J_4+J_5$, and by the Wigner--Eckart theorem for
a scalar product of tensor operators on two subsystems~\cite[(7.1.6)]{Edmonds},
\[
  \langle b-1,c|X|b,c\rangle=(-1)^{b+c+s}\begin{Bmatrix}s&c&b-1\\1&b&c\end{Bmatrix}
  \langle b-1\|J_2\times J_3\|b\rangle\,\langle c\|L\|c\rangle .
\]
The $6j$ symbol equals $\{s\ c\ b;\,1\ b-1\ c\}$, whose sign is $(-1)^{s+b+c}$
\cite[Table~5]{Edmonds}; the two $c$-dependent signs cancel, $\langle c\|L\|c\rangle>0$,
and the remaining factor does not involve $c$. So the phase of an $X$-link depends on
$b$ alone, and likewise the phase of a $D$-link depends on $c$ alone. A diagonal
unitary depending on $(b,c)$ then makes every link positive.
\end{proof}

The identity \eqref{eq:heron} is where the Heron factors of the $6j$ coefficient
split: the $s$-dependence of every link is a function of the diagonal $d$ times a
function of the antidiagonal $u$. It is the source of the congruence in
\cref{thm:heron-stark}. The sites carry a natural two-colouring, the parity of
$u-s$, and every link joins the two colours, so the matrix is bipartite and its
nonzero spectrum is $\pm$ the nonzero singular values of the block $B$ from one colour
to the other. \Cref{fig:grid} draws a small example.

\begin{schemfigure}
\begin{tikzpicture}[scale=0.78]
  \fill[esPlotGrid!70] (-0.45,-0.45) rectangle (6.45,6.45);
  \draw[esPlotGrey, densely dotted] (-0.45,-0.45) rectangle (6.45,6.45);
  \fill[esPlotBlue!12] (-0.45,0.85) -- (5.15,6.45) -- (6.45,6.45) -- (6.45,5.15) -- (0.85,-0.45) -- (-0.45,-0.45) -- cycle;
  \draw[esPlotBlue!70, dashed] (-0.45,0.55) -- (5.45,6.45);
  \draw[esPlotBlue!70, dashed] (0.55,-0.45) -- (6.45,5.45);
  \foreach \a/\b/\c/\d in {0/1/1/1,1/1/2/1,1/2/2/2,2/2/3/2,2/3/3/3,3/3/4/3,3/4/4/4,4/4/5/4,4/5/5/5,5/5/6/5,5/6/6/6} \draw[esPlotBlue, line width=1.1pt, shorten >=4pt, shorten <=4pt] (\a,\b) -- (\c,\d);
  \foreach \a/\b/\c/\d in {1/0/1/1,1/1/1/2,2/1/2/2,2/2/2/3,3/2/3/3,3/3/3/4,4/3/4/4,4/4/4/5,5/4/5/5,5/5/5/6,6/5/6/6} \draw[esPlotAmber, line width=1.1pt, shorten >=4pt, shorten <=4pt] (\a,\b) -- (\c,\d);
  \foreach \a/\b in {0/1,1/0,1/2,2/1,2/3,3/2,3/4,4/3,4/5,5/4,5/6,6/5} \fill[esPlotInk] (\a,\b) circle (3.4pt);
  \foreach \a/\b in {1/1,2/2,3/3,4/4,5/5,6/6} \draw[esPlotInk, fill=white, line width=0.9pt] (\a,\b) circle (3.4pt);
  \draw[esPlotInk, line width=0.8pt] (-0.15,-0.15) -- (0.15,0.15) (-0.15,0.15) -- (0.15,-0.15);
  \draw[->] (-0.45,-0.9) -- (7.0,-0.9) node[right] {$b$};
  \draw[->] (-0.9,-0.45) -- (-0.9,7.0) node[above] {$c$};
  \node[font=\small, text=esPlotBlue!85!black, anchor=west] at (4.6,7.0) {wall $d=-s$};
  \node[font=\small, text=esPlotBlue!85!black, anchor=west] at (6.6,5.3) {wall $d=s$};
  \node[font=\small, anchor=west, align=left] at (8.0,2.3)
     {\textcolor{esPlotBlue!85!black}{horizontal links: $X=J_1\cdot(J_2\times J_3)$,}\\
      \textcolor{esPlotBlue!85!black}{\quad Stark link $h_1(b)$ times Heron factors}\\[2pt]
      \textcolor{esPlotAmber!85!black}{vertical links: $D=J_1\cdot(J_4\times J_5)$,}\\
      \textcolor{esPlotAmber!85!black}{\quad Stark link $h_2(c)$ times Heron factors}\\[2pt]
      filled and open: the two colours,\\ the parity of $u-s$};
\end{tikzpicture}
\end{schemfigure}
\esfiglabel{fig:grid}%
\schemcap{FIG.\ \theesfig}{The recoupling grid for $s=j_1=1$ and two pairs whose spins
range over $0,\dots,6$; the cross marks the corner site $(0,0)$, excluded because
$b+c<s$. The states fill the rectangle of the two pair spins $(b,c)$, cut
to the strip $|b-c|\le s$ and above the corner $b+c\ge s$. The first triple product moves horizontally, the second
vertically, and every link joins the two colours. For $s=\frac12$ the strip is a
single zigzag path: the chain of \cref{sec:closing}.}

\begin{proposition}{Symmetry at equal weights}{S4}
At $t=1$, with $T=J_1\cdot(J_2\times J_3+J_4\times J_5)$, the moduli of the
spectrum of $T$ are invariant under every permutation of $j_2,j_3,j_4,j_5$. On the
invariant space $T=-\sum_{2\le a<b<c\le5}J_a\cdot(J_b\times J_c)$ and
$T=-i\,[J_2\cdot J_3-J_4\cdot J_5,\;J_3\cdot J_4-J_2\cdot J_5]$.
\end{proposition}

\begin{proof}
On invariant vectors $T=-V\cdot L$ with $V=J_2\times J_3+J_4\times J_5$ and
$L=J_2+\cdots+J_5$, because $J_1$ commutes with $V$; and $(J_2\times J_3)\cdot(J_2+J_3)=0$
since $J_2\cdot(J_2\times J_3)=iJ_2\cdot J_3=-J_3\cdot(J_2\times J_3)$. What remains are
the four triple products of the other legs, a cyclically symmetric expression; the
commutator form follows from $J_a\cdot(J_b\times J_c)=\frac i4[(J_a+J_b)^2,(J_b+J_c)^2]$.
The transposition of legs $2,3$ maps $T$ to $-X+D$, which the gauge $(-1)^b$ maps back
to $X+D$. A transposition and the four-cycle generate the symmetric group.
\end{proof}
